\documentclass[11pt,a4paper]{article}
\usepackage[utf8]{inputenc}
\usepackage[spanish]{babel}
\usepackage{amsmath, amssymb}
\usepackage{graphicx}
\usepackage[margin=2.5cm]{geometry}
\usepackage{fancyhdr}
\usepackage{tcolorbox}
\usepackage{enumitem}

% Configuración de colores corporativos UCB
\definecolor{UCBblue}{RGB}{0, 51, 153}
\definecolor{UCBgray}{RGB}{100, 100, 100}

% Configuración de encabezado y pie de página
\pagestyle{fancy}
\fancyhf{}
\fancyhead[L]{\textcolor{UCBblue}{\textbf{Ingeniería Mecatrónica - UCB Tarija}}}
\fancyhead[R]{\textcolor{UCBgray}{Taller de Nivelación - Módulo 1}}
\fancyfoot[C]{\thepage}
\renewcommand{\headrulewidth}{0.4pt}
\renewcommand{\headrule}{\hbox to\headwidth{\color{UCBblue}\leaders\hrule height \headrulewidth\hfill}}

\begin{document}

\begin{center}
    {\Large \textbf{\textcolor{UCBblue}{Guía de Modelado Analítico: Sistematización de Redes de Potencia}}}\\
    \vspace{0.3cm}
    {\large \textbf{Fase CDIO: Concebir (Modelado Físico-Matemático)}}\\
    \vspace{0.5cm}
    \textbf{Estudiante:} \rule{8cm}{0.4pt} \quad \textbf{Fecha:} \rule{3cm}{0.4pt}
\end{center}

\vspace{0.5cm}

\section*{1. Fundamentación Física: La Red de Potencia}
El diseño de sistemas mecatrónicos requiere el dimensionamiento preciso de las etapas de potencia (drivers, cableado, protecciones). Para modelar matemáticamente este comportamiento, recurrimos a analogías fundamentales:
\begin{itemize}
    \item \textbf{Voltaje (V):} Representa el gradiente de potencial o fuerza electromotriz. Es análogo a la presión generada por una bomba en un sistema hidráulico.
    \item \textbf{Corriente (I):} Cuantifica el flujo de carga a través de un conductor (Amperios), análogo al caudal volumétrico.
    \item \textbf{Resistencia (R):} Cuantifica la oposición geométrica y material al flujo de carga (Ohmios).
    \item \textbf{Malla (Lazo Independiente):} Trayectoria cerrada elemental dentro del circuito que no contiene otras trayectorias en su interior.
\end{itemize}

\section*{2. Sistematización Algorítmica: Formulación Matricial (LVK)}
La resolución computacional de circuitos mediante herramientas científicas (Python/MATLAB) requiere transponer la topología física a una estructura de datos matricial canónica: $\mathbf{R}\mathbf{I} = \mathbf{V}$. 

Para ensamblar la \textbf{Matriz de Resistencias de Lazo ($\mathbf{R}$)}, aplique el siguiente procedimiento sistemático:

\begin{tcolorbox}[colback=blue!5!white,colframe=UCBblue,title=\textbf{Procedimiento de Ensamblaje Matricial}]
\textbf{Paso 1: Convención de Estado} \\
Asigne convencionalmente una corriente de lazo en sentido horario ($\circlearrowright$) a cada malla independiente. Denomínelas $I_1, I_2, \dots, I_N$.

\vspace{0.2cm}
\textbf{Paso 2: Elementos de la Diagonal Principal ($R_{kk}$)} \\
La impedancia propia de la malla $k$ (elemento $R_{kk}$) equivale a la suma algebraica de \textbf{todas} las resistencias que conforman dicha malla. \\
\textit{Condición de viabilidad: En redes pasivas, $R_{kk}$ es estrictamente positivo ($R_{kk} > 0$).}

\vspace{0.2cm}
\textbf{Paso 3: Elementos Extradiagonales ($R_{kj}$)} \\
La impedancia mutua entre la malla $k$ y la malla $j$ (elemento $R_{kj}$) corresponde al valor de la resistencia compartida en la frontera topológica de ambas mallas, \textbf{precedida por un signo negativo}. \\
\textit{Condición de viabilidad: La matriz resultante debe ser estructuralmente simétrica ($R_{kj} = R_{jk}$).}
\end{tcolorbox}

\section*{3. Condición de Excitación (Vector $\mathbf{V}$)}
Para formular el vector columna de fuentes independientes $\mathbf{V}$, evalúe el recorrido de la corriente de malla a través de los elementos activos (baterías/fuentes):
\begin{itemize}
    \item Si el recorrido horario ingresa por el terminal \textbf{negativo} de la fuente, se considera una elevación de potencial y el valor se ingresa como \textbf{positivo} en el vector.
    \item Si ingresa por el terminal \textbf{positivo}, representa una caída de potencial y se ingresa como \textbf{negativo}.
\end{itemize}

\newpage

\section*{4. Desarrollo Analítico: Misión de Ingeniería}

\begin{figure}[h!]
    \centering
    % Placeholder para la imagen de LTSpice
    \framebox[14cm][c]{\rule{0pt}{6cm} \textit{[Espacio reservado: Inserte aquí la exportación del esquemático desde LTSpice]}}
    \caption{Topología coplanar de la red de distribución de potencia mecatrónica (3 Mallas independientes).}
\end{figure}

Considere la topología superior con los siguientes parámetros de hardware:
$$R_1=10\Omega, \quad R_2=22\Omega, \quad R_3=15\Omega, \quad R_4=33\Omega, \quad R_5=47\Omega, \quad R_6=12\Omega$$
$$V_{s1} = 24\text{ V (Lazo 1)}, \quad V_{s2} = 0\text{ V (Lazo 2)}, \quad V_{s3} = -12\text{ V (Lazo 3)}$$

\textbf{A.} Ejecute el procedimiento de sistematización algorítmica para estructurar analíticamente el sistema canónico de ecuaciones $\mathbf{R}\mathbf{I} = \mathbf{V}$. Complete el tensor:

\vspace{0.5cm}

$$
\begin{bmatrix}
 (R_1 + \dots \dots + R_3) & -R_2 & -R_3 \\
 \dots \dots \dots \dots \dots & (\dots \dots + R_4 + \dots \dots) & -R_5 \\
 \dots \dots \dots \dots \dots & \dots \dots \dots \dots \dots & (R_3 + \dots \dots + R_6)
\end{bmatrix}
\begin{bmatrix}
 I_1 \\
 I_2 \\
 I_3
\end{bmatrix}
=
\begin{bmatrix}
 24.0 \\
 \dots \dots \\
 \dots \dots
\end{bmatrix}
$$

\vspace{0.5cm}

\textbf{B.} Reemplace los parámetros por sus valores nominales y evalúe el determinante principal de la matriz ($\det(\mathbf{R})$). Justifique físicamente por qué $\det(\mathbf{R}) \neq 0$ es un requerimiento crítico para descartar condiciones de cortocircuito o inconsistencias topológicas.

\vspace{0.3cm}
\begin{tcolorbox}[colback=white,colframe=UCBgray,height=6.5cm,valign=top]
\textit{Espacio para el desarrollo del cálculo (Regla de Sarrus o Cofactores) y justificación física:}
\end{tcolorbox}

\section*{5. Transición al Entorno Computacional (Gemelo Digital)}
Una vez certificado su modelo matemático analítico, proceda a la fase de Implementación:
\begin{enumerate}[label=\textbf{\arabic*.}]
    \item Inicialice su entorno de desarrollo local (Jupyter Notebooks).
    \item Abra el archivo de trabajo asignado: \texttt{Notebook\_01\_Mallas.ipynb}.
    \item Declare la matriz $\mathbf{R}$ utilizando arreglos bidimensionales en \texttt{NumPy}.
    \item Ejecute la celda de pruebas unitarias (\texttt{assert}). Si el código valida la simetría y positividad de la diagonal, su modelo es computacionalmente robusto.
\end{enumerate}

\end{document}
